\documentclass[11pt, a4paper]{article}
\usepackage[utf8]{inputenc}
\usepackage[T1]{fontenc}
\usepackage[french]{babel}
\usepackage{geometry}
\geometry{margin=2.5cm}
\usepackage{amsmath}
\usepackage{amssymb}
\usepackage{hyperref}
\usepackage{listings}
\usepackage{xcolor}
\usepackage{tcolorbox}
\usepackage{graphicx}

\definecolor{codegray}{gray}{0.95}
\lstset{
    backgroundcolor=\color{codegray},
    basicstyle=\ttfamily\small,
    breaklines=true,
    frame=single
}

\title{Rapport Technique Exhaustif : Mesh Comparison Viewer}
\author{Système d'analyse Antigravity}
\date{\today}

\begin{document}

\maketitle

\tableofcontents
\newpage

\section{Introduction et Contexte}
Ce projet a été mis en place pour pallier le besoin de valider quantitativement et qualitativement les reconstructions de surfaces d'objets manipulés à partir de données de contact discrètes (issues du dataset \textit{grasp-dataset-gen}). 

L'application fournit un pipeline de calcul en temps réel combinant la reconstruction implicite par transformée de Fourier, la voxélisation unifiée (grille standardisée), ainsi que l'extraction de métriques volumiques.

\section{Architecture Logicielle}
L'environnement a été conçu de façon modulaire pour séparer le chargement, le calcul algorithmique pur, et la couche de rendu graphique.

\subsection{Organisation des sous-modules (\texttt{core/})}
\begin{itemize}
    \item \textbf{\texttt{data\_loader.py}} :
        Il s'interface avec le système de fichiers du générateur de dataset. Il normalise automatiquement les modèles 3D (centrage sur le centroïde et mise à l'échelle à un rayon unifié de 8cm) pour assurer la robustesse des étapes de voxélisation suivantes.
    \item \textbf{\texttt{poisson.py}} :
        Implémente le solveur FFT. Il encapsule la transformation des points spatiaux en coordonnées de grille, le calcul de la divergence vectorielle $\nabla \cdot \vec{V}$ par gradients finis Numpy, la résolution matricielle fréquentielle, et l'isosurfaçage final via \textit{skimage.measure.marching\_cubes}.
    \item \textbf{\texttt{voxelizer.py}} :
        Ce module opère sur une grille discrète immuable : un cube virtuel de $20$ cm centrée en $(0,0,0)$ découpée en $128 \times 128 \times 128$ sous-cubes (pitch $\approx 1.56$ mm). Il traite à la fois la conversion booléenne de maillage watertight et le ré-échantillonnage par interpolation trilinéaire (\texttt{RegularGridInterpolator}) des champs scalaires implicites.
    \item \textbf{\texttt{visualization.py}} :
        Convertit les structures de données complexes (Trimesh, tableaux Numpy 3D) en traces \textit{Plotly}. Il gère l'habillage "Premium Dark" : shaders pour les maillages, scatter volumique optimisé pour les nuages de voxels, et cartographie de chaleur (\textit{heatmap}) interactives pour les coupes 2D.
\end{itemize}

\section{Procédés Algorithmiques en Profondeur}

\subsection{Reconstruction par Screened Poisson (Approximation FFT)}
L'objectif est de trouver la fonction indicatrice $\chi$ d'une surface fermée dont le gradient matche au mieux le champ vectoriel des normales $V$ donné par les points de contact.
L'approche par FFT suit ces étapes rigoureuses :
\begin{enumerate}
    \item \textbf{Convolution Spatiale} : Chaque point de contact $p_i$ et sa normale $\vec{n}_i$ sont "projetés" sur une grille régulière $N \times N \times N$. La contribution à une cellule $x_{ijk}$ suit une Gaussienne : $W(d) = \exp(-d^2 / 2\sigma^2)$.
    \item \textbf{Passage au Champ Scalaire} : On calcule la divergence du champ obtenu : $f = \nabla \cdot \vec{V}$.
    \item \textbf{Résolution dans le Domaine Fréquentiel} : Dans l'espace de Fourier, l'équation de Poisson $\Delta\chi = f$ devient algébrique : 
    $ \mathcal{F}(\chi)(\vec{k}) = \frac{\mathcal{F}(f)(\vec{k})}{- \|\vec{k}\|^2} $
    On applique la transformée inverse $\mathcal{F}^{-1}$ pour récupérer $\chi$ dans l'espace réel.
    \item \textbf{Seuillage d'Iso-surface} : La valeur $\chi$ est échantillonnée aux coordonnées exactes des points de contact initiaux. La moyenne arithmétique de ces valeurs sert de "niveau zéro" (isovaleur) pour reconstruire le maillage via les Cubes Marchants.
\end{enumerate}

\subsection{Unification de la Grille de Voxelisation}
Pour comparer le Maillage de Référence (Ground Truth) et le champ SDF continu généré par Poisson, nous forçons les deux entités dans le même espace tensoriel booléen.
\begin{itemize}
    \item \textbf{Traitement GT} : La surface est transformée en sous-grille via l'algorithme de scan-conversion de \texttt{trimesh}, puis un algorithme de \textit{Flood Fill} 3D scelle le volume intérieur.
    \item \textbf{Traitement SDF} : Le cube standard query la valeur interpolée de $\chi$ à chaque centre de voxel. Si $\chi(x,y,z) \geq \text{Isovaleur} + \text{Offset}$, le voxel est déclaré "intérieur" (Occupancy=True).
\end{itemize}

\section{Analyse des Paramètres et Sensibilité}

\begin{tcolorbox}[colback=gray!5!white,colframe=gray!75!black,title=Comportement des leviers de réglage]
\begin{itemize}
    \item \textbf{Poisson Resolution} : Contrôle le pas d'échantillonnage de la FFT. \textit{Conséquence} : Une valeur basse lisse le modèle et comble artificiellement les trous entre les doigts. Une valeur haute peut provoquer du "bruit" si les points de contact sont trop espacés.
    \item \textbf{Sigma Factor} : Contrôle la largeur du noyau Gaussien. \textit{Conséquence} : Si $\sigma$ est faible, la normale n'influence que sa cellule immédiate (reconstruction pointue et locale). S'il est élevé, la normale influence un large voisinage (effet blob global).
    \item \textbf{SDF Threshold Offset} : Modifie le biais de la classification. \textit{Conséquence} : Réduit ou augmente mathématiquement le rayon volumique de l'objet généré, ce qui permet d'affiner manuellement la "coque" de saisie si l'isosurfaçage moyen a sous-estimé la pénétration.
\end{itemize}
\end{tcolorbox}

\section{Moteur de Métriques et Interprétation}

L'outil calcule une matrice de confusion voxel-à-voxel où Vrai Positif (VP) = les deux versions s'accordent sur l'occupation.

\begin{enumerate}
    \item \textbf{IoU (Intersection over Union) / Jaccard Index} :
    \[ \text{IoU} = \frac{|GT \cap P|}{|GT \cup P|} \]
    Mesure absolue de chevauchement. Varie entre 0 (aucune intersection) et 1 (identique). C'est le standard le plus sévère car il pénalise les débordements.
    \item \textbf{Dice Coefficient (F1-Score 3D)} :
    \[ \text{Dice} = \frac{2 \cdot |GT \cap P|}{|GT| + |P|} \]
    Donne plus de poids à l'intersection. Idéal pour visualiser si la masse globale du volume est capturée.
    \item \textbf{Precision vs Recall} :
    \begin{itemize}
        \item \textbf{Precision} = $VP / P$. Si elle est faible, l'algorithme de Poisson a "inventé" du volume qui n'existe pas en vrai.
        \item \textbf{Recall} = $VP / GT$. Si il est faible, cela signifie que la reconstruction a manqué des parties entières de l'objet réel (dû au manque de points de contact sur ces zones).
    \end{itemize}
\end{enumerate}

\end{document}
